% Writing and reference shorthands retained from the author's archived papers.
% Mathematical notation belongs in math_commands.tex.

% Figure-panel prose.
\newcommand{\figleft}{\emph{(Left)}}
\newcommand{\figcenter}{\emph{(Center)}}
\newcommand{\figright}{\emph{(Right)}}
\newcommand{\figtop}{\emph{(Top)}}
\newcommand{\figbottom}{\emph{(Bottom)}}
\newcommand{\captiona}{\emph{(a)}}
\newcommand{\captionb}{\emph{(b)}}
\newcommand{\captionc}{\emph{(c)}}
\newcommand{\captiond}{\emph{(d)}}
\newcommand{\newterm}[1]{\textbf{#1}}

% Compatibility wrappers around cleveref. New prose may also call \cref and
% \Cref directly.
\newcommand{\figref}[1]{\cref{#1}}
\newcommand{\Figref}[1]{\Cref{#1}}
\newcommand{\twofigref}[2]{\cref{#1,#2}}
\newcommand{\quadfigref}[4]{\cref{#1,#2,#3,#4}}
\newcommand{\secref}[1]{\cref{#1}}
\newcommand{\Secref}[1]{\Cref{#1}}
\newcommand{\twosecrefs}[2]{\cref{#1,#2}}
\newcommand{\secrefs}[3]{\cref{#1,#2,#3}}
\newcommand{\algref}[1]{\cref{#1}}
\newcommand{\Algref}[1]{\Cref{#1}}
\newcommand{\twoalgref}[2]{\cref{#1,#2}}
\newcommand{\Twoalgref}[2]{\Cref{#1,#2}}
\newcommand{\plaineqref}[1]{\ref{#1}}

% Drafting and table helpers.
\newcommand{\red}[1]{\textcolor{red}{#1}}
\newcommand{\blue}[1]{\textcolor{blue}{#1}}
% jmlr2e defines this proof-ending symbol; retain a fallback for other styles.
\providecommand{\BlackBox}{\rule{1.5ex}{1.5ex}}
\newcommand{\cmark}{\ding{51}}
\newcommand{\xmark}{\ding{55}}

% Pseudocode assignment operators.
\newcommand{\pluseq}{\mathrel{{+}{=}}}
\newcommand{\minuseq}{\mathrel{{-}{=}}}
\newcommand{\prodeq}{\mathrel{{*}{=}}}
